\documentclass[10pt]{article} % use larger type; default would be 10pt
\usepackage[utf8]{inputenc} % set input encoding (not needed with XeLaTeX)

\usepackage{geometry} % to change the page dimensions
\geometry{letterpaper}
\usepackage{graphicx} % support the \includegraphics command and options
\setlength\voffset{-0.2in}

\title{\includegraphics[width=1.55in]{ruddock_crest.png} \\ Constitution}
\author{Ruddock House}
%\date{} % Activate to display a given date or no date (if empty),

\begin{document}

\maketitle

\section{Name and Membership}
\subsection{Name}
The name of the organization shall be Ruddock House.
\subsection{Membership}
\begin{enumerate}
\item The types and qualifications of membership are:
\begin{enumerate}
\item a Full Member, who is an undergraduate of the California Institute of Technology and is picked as such at the end of Rotation Week, or is elected as such by a 2/3 majority of those voting at a House meeting;
\item a Social Member, who is an undergraduate at the California Institute of Technology, is not residing in Ruddock House or a designated off-campus alley and has been elected as such, following a week of consideration, by a simple majority of those voting, when those voting constitute at least 1/3 of the House Membership;
\item an Associate Member, who was an undergraduate of the California Institute of Technology and was a Full Member or Social Member of Ruddock House, or who is any other person associated with Ruddock House and so designated by the Executive Committee;
\item the Resident Associates, who shall be considered to be Full Members.
\end{enumerate}
\item Full and social members who discontinue paying their house dues while they are registered students at the California Institute of Technology must be re-elected at a house meeting as provided for in Section 1.2, part 1a.
\item Full members can voluntarily demote themselves to social membership or cancel their membership. Social members can voluntarily cancel their membership.
\end{enumerate}
\subsection{Duties and Privileges}
\begin{enumerate}
\item The duties and privileges of Full Members apply immediately upon attaining such status. These duties and privileges are:
\begin{enumerate}
\item to attend all regular and special House Meetings;
\item to participate actively in all House functions;
\item to pay dues;
\item to have room preference;
\item to vote on all matters put before the House;
\item to abide by the rules and regulations made by the Executive Committee, the Upper Class Committee, and the Committee on Student Houses.
\end{enumerate}
\item The duties and privileges of Social Members apply immediately upon ratification by The House and are identical to Full Members, with the exception of:
\begin{enumerate}
\item the right to have room preference;
\item the right to vote on all matters put before The House;
\end{enumerate}
\item The duties and privileges of associate members:
\begin{enumerate}
\item do not include a vote in House matters;
\item shall be determined by the Executive Committee in other extents.
\end{enumerate}
\end{enumerate}
\section{Officers}
\subsection{General Qualifications}
Each officer shall be a Full Member of the House when elected and during his term of office. If, during his term of office, his membership changes, he immediately forfeits his office.
\subsection{Qualifications and Duties of Elected Officers}
\begin{enumerate}
\item The President shall reside in Ruddock House and shall be a junior or senior for the latter portion of his term in office. He shall represent the House on the Interhouse Committee, preside over House Meetings, be chairman of the Executive Committee, and be an ex-officio member of all House committees, and shall be responsible for selecting Full Members to attend freshmen picks.
\item The Vice-president shall meet the same qualifications as the President. He shall assume the duties of the President in the event of the President’s absence. He shall be chairman of the Upper Class Committee. He shall be responsible for maintaining the external relations of the House, for bringing guests to the House, and shall act as official host for all House guests. He shall be responsible for the entire freshman selection except as noted above.
\item The Secretary shall reside in Ruddock House. He shall record minutes, including all motions made and the votes on these by all officers present, of all House, Executive Committee, and Upper Class Committee meetings; and post publicly a copy of the minutes of all open meetings. He shall have custody of all House records not specifically delegated to the custody of another officer, and in general shall perform all the duties of a corresponding and recording secretary. He shall be responsible for room assignments.
\item The Treasurer shall reside in Ruddock House. He shall be responsible to the House for all receipts and expenditures, shall maintain an adequate book-keeping system, shall submit a report to the Executive Committee at each regularly scheduled meeting, shall propose a budget to the Executive Committee for each term, shall submit to the House a written report of the finances at the end of each term, and shall maintain house property.
\item The Librarian shall reside in Ruddock House and shall be a sophomore for the latter portion of his term in office. He shall be responsible for procuring magazines and newspapers, and for maintaining the House Library, including its records. He shall publish a House list at the beginning of each term. He shall be responsible for freshman orientation and organizing the freshman members of the House for the performance of those tasks assigned to them.
\item The office of Social Chairman shall be held by not more than five members of the House and will be responsible for organizing all dances, parties and other social functions, and for publishing a calendar of such events.
\item The office of Athletic Manager shall be held by not more than three members of the House. The Athletic Team shall be responsible for all athletic equipment used by the House, for the participation of the House in Interhouse and intrahouse athletics, and in general shall organize and promote the athletic activities of the House.
\item The order of succession shall be President, Vice-president, Secretary, Treasurer, Librarian, Social Chairman, and Athletic Manager.
\item In the event that the President is unable to complete his term of office, the Vice-President shall assume all responsibilities of the President, and a new Vice-president shall be elected pursuant to nomination and election guidelines.
\end{enumerate}
\subsection{Duties and Qualifications of Appointed Officers}
\begin{enumerate}
\item One Upper Class Committeeman shall reside in each Ruddock House Alley and two shall reside off-campus\footnote{We denote an off-campus member of Ruddock House as a member who does not physically reside in the house.}, except in the case where no suitable volunteers can be found living off-campus, in which case the Off-Campus Upper Class Commiteemen may live on-campus. An Upper Class Commiteeman shall have been both enrolled at Caltech for at least five academic terms and a Full House Member for at least three academic terms prior to his term of office. He shall be available to counsel the occupants of the alley in which he resides or, in the case of an Off-Campus Upper Class Commiteemen, any off-campus member. He shall be responsible for maintaining order and enforcing the rules and policies of the House. He shall not concurrently hold a position on the Executive Committee and be an Upper Class Committeeman.
\item The Historian shall be responsible for all recording and preservation of House records and traditions, including House scrapbook, House history, and House roll, and for submission of the House article in the Institute yearbook.
\item The Head Waiter shall be responsible for maintaining proper decorum and conduct in the dining room, and shall direct the student waiters.
\item The House BFD Editors shall be responsible for publishing the House paper.
\item The House O’Domhnaill’s Supplier shall maintain the House convenience store.
\end{enumerate}
\subsection{Nominations, Elections, and Appointments}
\begin{enumerate}
\item Nominations for ARC representative, CRC representative, two BoC representatives, Athletic Manager, Social Chairman, Librarian, Treasurer, Secretary, Vice-president, and President, in that order, shall be made at a meeting of the House held for that purpose between the beginning of the second term and the third Friday of second term. Nominations shall remain open until the election of that office.
\item Election of all officers shall be carried out at a House Meeting held for that purpose within one week after the nomination meeting. The order of election shall be in the following order: President, Vice-president, Secretary, Treasurer, Librarian, Social Chairman, Athletic Manager, two BoC representatives, CRC representative, ARC representative. The election for each office shall be held before nominations for successive offices are closed. Election shall require a majority of votes cast from the total voting membership, not including blank ballots, and they shall be made by secret ballot. Elections shall be managed by an Election Committee in accordance with the Procedural Amendment.
\item For the election of the BoC representatives, each voting member shall vote for two candidates per round. For each other office, each voting member shall vote for one candidate per round.
\item All officers shall be installed at a meeting of the House held not more than one week from their election. The oath of office shall be administered by the retiring president. Each officer shall assume the duties of office immediately upon being installed.
\item Sign-up lists for appointed offices shall be posted between the first day of Third Term the third Friday of Third Term and sign-up lists shall be removed and appointments made by the Executive Committee two weeks thereafter. The term of service for appointed offices shall be the academic year directly following their appointment by the Executive Committee.
\item Except as stated in Section 2.2, Paragraph 9, if a vacancy in an elected or appointed office occurs, the Executive Committee, or the remaining part thereof, shall select a qualified member of the House to fill the vacancy until the next election or appointment process.
\end{enumerate}
\subsection{Recall}
\begin{enumerate}
\item Any elected officer may be recalled by a majority vote of the voting membership of the House. 
\item The procedure shall be as follows:
\begin{enumerate}
\item Petitions, containing reasons and purpose, must be submitted in writing, endorsed by 1/3 of the total Full Membership, at a House meeting held for that purpose;
\item The petitions must be posted immediately following the presentation and remain posted until voting is closed;
\item Voting shall follow the specified procedures for Constitutional amendment elections.
\end{enumerate}
\end{enumerate}
\section{The Executive Body}
\subsection{Name}
The executive body of the House shall be known as the Executive Committee.
\subsection{Membership}
The Executive Committee shall consist of the President, who shall act as chairman, the Vice-president, Secretary, Treasurer, Librarian, Social Chairman, and Athletic Manager.
\subsection{Voting}
All members of the Executive Committee shall vote. Each office shall receive one vote, whether filled by one or more people.
\subsection{Powers}
\begin{enumerate}
\item The Executive Committee shall formulate House rules and policies, and assume all duties not otherwise delegated.
\item It shall appoint all offices described under Section 2.3 of this Constitution.
\item It shall designate associate members.
\item It shall have the power of dismissal over all appointed officers, other than Upper Class Committeemen.
\end{enumerate}
\subsection{Meetings}
\begin{enumerate}
\item Meetings may be called by the President or any two members.
\item A quorum shall consist of four of the seven offices represented in a voting capacity.
\item There shall be regularly scheduled meetings, at least biweekly, the day to be determined by the Executive Committee.
\item Closed meetings, attended only by the Executive Committee members, shall be held when discussing individuals and other sensitive issues, at the discretion of the Executive Committee.
\item Upon petition of 1/3 of the Full Membership of the House, a meeting of the Executive Committee shall be held within two weeks to consider any proposal so petitioned.
\end{enumerate}
\section{The Judiciary Body}
\subsection{Name}
The judiciary body of the House shall be known as the Upper Class Committee.
\subsection{Membership}
The Upper Class Committee shall consist of the Upper Class Committeemen.
\subsection{Powers and Duties of the Committee}
The Upper Class Committee:
\begin{enumerate}
\item shall review and act on all violations of House rules and policy, by reprimand, punishment, or other appropriate action.
\item shall be responsible for the administration and tabulation of recall and amendment elections.
\item shall be autonomous in all matter of judicial policy not otherwise specified in this Constitution.
\end{enumerate}
\subsection{Meetings}
\begin{enumerate}
\item Meetings may be called by the Vice-president, Resident Associate, or any two Upper Class Committeemen.
\item The Vice President shall preside over Upper Class Committee meetings and shall vote only in the case of a tie.
\item A quorum shall consist of at least 2/3 of the Upper Class Committeemen, provided that when the Committee considers the case of an individual person, the committeeman of the alley in which the person lives, and the committeeman of the alley in which the infraction occurred, must be present.
\end{enumerate}
\section{House Meetings}
\subsection{Meetings}
\begin{enumerate}
\item The House shall meet:
\begin{enumerate}
\item for the nomination, election, and recall of officers;
\item when a meeting is called by the Executive Committee;
\item at other times, on petition by 1/5 of the Full Membership.
\end{enumerate}
\item Notice of the meeting shall be published at least three days prior to each House meeting in a manner expected to reach the entire membership of the House.
\begin{enumerate}
\item The notice shall contain the date, time, and place of the meeting and an agenda of all issues to be discussed and voted upon at that meeting.
\item No substantive issues not on a meeting’s posted agenda may be voted upon at that meeting. Discussion of any issue and voting on non-substantive issues are not constrained.
\end{enumerate}
\end{enumerate}
\subsection{Quorum}
\begin{enumerate}
\item A quorum shall consist of the larger of:
\begin{enumerate}
\item 1/3 of the Full-Membership of the House, or
\item 80\% of the average attendance of House meetings held during the three hundred sixty-five (365) days prior to the announced date of the meeting regardless of whether those meetings established quorum.
\end{enumerate}
\item If quorum is not established, no substantive matters can be decided.
\end{enumerate}
\subsection{Procedure}
In case of a dispute on parliamentary procedure in meetings of the House or any of its committees, Robert’s Rules of Order (revised edition) shall provide the rule whenever the Constitution does not.
\section{Finances}
\subsection{Dues}
\begin{enumerate}
\item The dues of the House shall be due and payable by each regular member at registration for each term. 
\item Married students shall be exempt from dues.
\item The dues for Social Members shall be set by a majority vote of the Executive Committee.
\item The dues of the Associate Members shall be determined by the Executive Committee.
\end{enumerate}
\subsection{Levies}
The dues may be changed, or addition assessments levied, only by a 1/2 vote of the Full Membership.
\subsection{Payment of Moneys}
\begin{enumerate}
\item Checks and other instruments for the payment of moneys shall be drawn in the name of the House and shall be signed either by the President or by the Treasurer.
\item The Treasurer may authorize House related expenditures up to \$250.
\item House related expenditures greater than \$250 must be approved by a simple majority vote of the Executive Committee.
\item Any expenditure may be placed to a House vote at the discretion of the Executive Committee. A simple majority of the House voting membership may approve any expenditure.
\end{enumerate}
\section{Room Hassles}
\subsection{Procedure}
\begin{enumerate}
\item The Secretary shall be responsible for running the Room Hassle in the spirit of fairness to maximize the occupancy of the house.
\item A Room Hassle shall be held:
\begin{enumerate}
\item once per year following the announcement of appointed offices to allocate rooms for the next academic year
\item whenever beds are available and there are Full Members seeking residence
\item whenever a partially filled room is sought by a member higher in the order given by section 7.1.4.
\end{enumerate}
\item A Room Hassle shall continue until either all the available beds have been filled or there are no longer any Full Members seeking residence.
\item Rooms shall be picked in the following order:
\begin{enumerate}
\item by the House President, then
\item by all other Full Members in order of
\begin{enumerate}
\item class (seniors then juniors...), then
\item intended number of occupants as a percentage of the room capacity, descending, then
\item office as specified in section 7.2.7, then
\item a random order prepared prior to the section 7.1.2a Room Hassle in a manner prescribed by
the Secretary.
\end{enumerate}
\item The Upper Class Committee shall arbitrate all conflicts and disputes.
\item The Secretary may consolidate unoccupied beds if doing so will increase the number of occupants of the house.
\end{enumerate}
\end{enumerate}
\subsection{Guarantees}
\begin{enumerate}
\item All Full Members shall have the opportunity to participate in any Room Hassle in person or by proxy.
\item All Full Members shall have the opportunity to designate their roommate(s) before another individual picks or is placed into the room. Thereafter, any unoccupied bed may be picked.
\item Contracted residents may exchange rooms and beds with each other at their own discretion. In the event of a dispute between roommates, the member with the higher pick shall have preference.
\item Room Hassles under section 7.1.2b or section 7.1.2c may not forcibly displace any participant from the house, and all non-participants shall retain their current rooms except as in 7.1.6. Members holding rooms sought under section 7.1.2c must participate in the hassle.
\item At any time during a Room Hassle, any participant may restart the process from his position in the order.
\item The President may have himself counted as a double occupant.
\item Entire rooms shall be guaranteed to the following officers:
\begin{enumerate}
\item In order of preference, House President, House Vice President, House Secretary, House Treasurer, Librarian, House Social Team (up to three members), Athletic Managers (up to two members), UCCs (one member per alley), IHC Chairman, ASCIT President, BoC Chair, House-elected BoC Representatives (up to two members), O’Domhnaill’s Suppliers (up to two members), House Head Waiter (one member), House BFD editors (one member),House Historians (one member)
\item for the entire academic year following the election or appointment provided the officer completes his term. If an officer quits his office mid-term, the room guarantee shall transfer to the succeeding officeholder for the remainder of the guarantee term.
\end{enumerate}
\item All Full Members who have picked into a room through a room hassle under 7.1.2a or 7.1.2b shall be guaranteed the same room upon their return from Study Abroad.
\end{enumerate}
\section{Amendments}
\subsection{Order}
Only Full Members may vote on Constitutional amendments.
\subsection{Necessary Vote}
An amendment to this Constitution passes if at least 2/3 of the votes cast are in favor of the amendment.
\subsection{Procedure}
\begin{enumerate}
\item Proposed amendments, containing reasons, purpose, and un-amended and amended language must be submitted in writing, endorsed by 1/5 of the Full Membership, at a House meeting.
\item The proposed amendments must be posted immediately following their presentation and remain posted until voting is closed.
\item Voting
\begin{enumerate}
\item Voting shall officially commence upon the legal presentation of the amendment and shall close at 11:59 PM on the seventh day following;
\item Voting shall be by secret ballot;
\item Votes of House Members shall not be examined prior to the closing of the election;
\item A roll of those members who have voted shall be available throughout the election.
\end{enumerate}
\end{enumerate}
\end{document}
